\chapter{Contrôle de risque pré-négociation et kill switches}\label{ch:ll:pre-trade-risk-and-kill-switches}

Le 1\up{er} août 2012, l'un des huit serveurs du routeur d'ordres d'un broker-dealer, Knight Capital Americas, portait encore un vieux
morceau de code que le nouveau déploiement aurait dû remplacer. Pour 212 ordres entrants, il envoyait des ordres enfants au marché «~sans
tenir compte du nombre d'exécutions d'actions~» déjà reçues, selon les termes de l'ordonnance du régulateur~: 4~millions d'exécutions sur 154
titres, plus de 397 millions d'actions, en environ 45 minutes, et une perte de 460~millions de dollars (le livre~7, chapitre~21, raconte l'histoire
du déploiement, et le livre~11, chapitre~27, la politique de risque qui en a découlé). Ce chapitre lit la même ordonnance comme
le ferait un ingénieur~: sous forme d'une liste d'exigences. Chaque paragraphe qui dit ce qui manquait à la société (un contrôle qui comparait les ordres
sortant du routeur avec ceux qui y étaient entrés, une limite de position qui comptait les ordres en cours, un seuil à l'échelle de la société qui
arrêtait automatiquement les ordres, un moyen de suspendre le routeur face à ses propres dysfonctionnements) est un contrôle à implémenter, sur le chemin de chaque
ordre, en quelques dizaines de nanosecondes. Le chapitre construit cette \omterm{def:ll:pre-trade-risk-and-kill-switches:gate}{barrière de risque}, puis rejoue à travers elle un emballement de la forme de 2012, avec
un contrôle à la fois, et constate que la plupart des contrôles auxquels une société penserait d'abord n'auraient pas suffi à l'arrêter.

\section{Le document d'exigences}

\begin{definition}[Barrière de risque, conception à défaillance fermée]\label{def:ll:pre-trade-risk-and-kill-switches:gate}
Une \emph{barrière de risque}\index{barrière de risque} (risk gate) est le composant placé sur le chemin de chaque ordre, entre le code qui décide de l'envoyer et
la \omterm{def:ll:order-gateway-and-order-management:gateway}{passerelle d'ordres}, qui le confronte aux limites de la société et le refuse si l'une d'elles est dépassée~; rien n'atteint la plateforme de négociation
sans l'avoir franchi. Une \emph{conception à défaillance fermée}\index{conception à défaillance fermée} (fail-closed design) est une conception dont les défaillances arrêtent l'activité qu'elle protège~:
une barrière de risque qui a perdu ses limites, ses prix de référence ou sa connaissance de la position de la société refuse les ordres plutôt
que de les laisser passer.
\end{definition}

L'ordonnance du régulateur (Release 34-70694) est précise sur ce qui manquait. La société disposait de contrôles avant que les ordres n'atteignent le
routeur, mais aucun après lui~: rien ne comparait «~les ordres qui sortaient de SMARS avec ceux qui y étaient entrés~», et il n'existait aucune procédure
«~pour arrêter les opérations de SMARS en réponse à sa propre activité aberrante~». Elle disposait d'un cap de prix, à 9.5~\% du meilleur prix lorsque l'ordre parent
est arrivé. Ce cap n'aidait pas tant que les cours bougeaient moins que cela et ne valait pas pour les ordres destinés à l'enchère
d'ouverture. Elle avait des limites de position pour certains groupes, mais «~ces limites ne tenaient pas compte de l'exposition de la société découlant
des ordres en cours~». Elle n'avait aucun seuil de capital à l'échelle de la société lié à des contrôles automatisés, et son suivi de position «~reposait entièrement sur
une surveillance humaine~». La règle à laquelle elle était soumise, la Market Access Rule, exige des contrôles qui empêchent l'entrée des ordres
qui dépassent des «~seuils de crédit ou de capital préétablis~» et des ordres erronés qui «~dépassent des paramètres de prix ou de taille
appropriés, au niveau de chaque ordre ou sur une courte période, ou qui indiquent des ordres en double~». La traduction technique
est une liste~: un tunnel de prix (price collar), des limites de taille et de notionnel par ordre, une limite de position qui compte les \omterm{def:ll:order-gateway-and-order-management:inflight}{ordres en vol}, un
seuil de capital à l'échelle de la société, une limite de débit, un contrôle des doublons et un kill switch que la barrière peut actionner elle-même. Le contrôle
de risque pré-négociation et le tunnel de prix sont définis dans le livre~11 (chapitre~27), qui en fixe la politique~; ce chapitre porte sur
la mise en œuvre rapide et correcte de ces contrôles.

\begin{dated}{2026-09}{Ce que les règles exigent des contrôles pré-négociation}\label{dat:ll:pre-trade-risk-and-kill-switches:rules}
Aux États-Unis, la Rule 15c3-5 du Securities Exchange Act (consulté en septembre 2026) impose à un courtier ou à un dealer ayant
accès au marché de disposer de contrôles qui empêchent l'entrée d'ordres dépassant des seuils de crédit ou de capital préétablis globalement,
pour chaque client et pour elle-même, et d'ordres erronés dépassant des paramètres de prix ou de taille, au niveau de chaque ordre ou sur une courte
période, ou indiquant des ordres en double. Dans l'Union européenne, le Commission Delegated Regulation (EU) 2017/589 exige qu'une
entreprise d'investissement de pouvoir annuler immédiatement tout ou partie de ses ordres non exécutés sur toutes les plateformes ou certaines d'entre elles (article 12)~; de disposer de
tunnels de prix, valeurs maximales par ordre, volumes maximaux par ordre et limites maximales de messages (article 15), et des
limiteurs d'exécution automatisés répétés~; et de générer des alertes en temps réel dans un délai de cinq secondes après l'événement concerné.
\end{dated}

La barrière s'insère dans le processus de trading, sur le thread du moteur de stratégie ou juste après lui, avant la \omterm{def:ll:order-gateway-and-order-management:gateway}{passerelle d'ordres}
(la \cref{fig:ll:pre-trade-risk-and-kill-switches:where})~: l'ordre du moteur ne devient le message de la passerelle que si la barrière
le valide. Tout le reste de la figure l'alimente. Les limites proviennent d'un service de risque, sous forme d'instantanés complets, avec un \omterm{def:ll:protocols-i-fix:session}{heartbeat}~; les
prix de référence viennent du \omterm{def:ll:the-order-book-builder:builder}{constructeur de carnet} (book builder)~; les exécutions et les annulations remontent de la passerelle, de sorte que la vue qu'a la barrière
de la position et des \omterm{def:ll:order-gateway-and-order-management:inflight}{ordres en vol} est celle de la passerelle~; et le kill switch peut être actionné par la barrière elle-même, par une personne
du desk de risque, ou par la plateforme de négociation.

\begin{omfigure}
\begin{tikzpicture}[font=\scriptsize,>=Stealth,box/.style={draw,rounded corners=2pt,minimum height=0.75cm,align=center,
  inner sep=3pt,fill=omDef!12}]
\node[box] (eng) at (0,0) {moteur\\ de stratégie};
\node[box,fill=omThm!18,minimum width=2.2cm,minimum height=1.1cm] (gate) at (3.3,0) {\textbf{barrière de risque}\\ 12 contrôles,\\ un seul masque};
\node[box] (gw) at (6.6,0) {passerelle\\ d'ordres};
\node[box,fill=gray!10] (venue) at (9.4,0) {plateforme};
\node[box,fill=omMeth!15] (lim) at (1.6,2.0) {service de risque~:\\ instantanés de limites,\\ heartbeats};
\node[box,fill=omMeth!15] (bb) at (5.0,2.0) {constructeur de carnet~:\\ prix de référence};
\node[box,fill=omMeth!15] (desk) at (3.3,-2.1) {desk de risque~:\\ kill switch};
\draw[->,thick] ([yshift=2mm]eng.east) -- node[above]{ordre} ([yshift=2mm]gate.west);
\draw[->,thick] (gate) -- node[above]{accepté} (gw);
\draw[->,thick] (gw) -- (venue);
\node[box,fill=gray!10] (mon) at (9.4,2.0) {moniteur\\ post-négociation};
\draw[->] (lim) -- (gate);
\draw[->] (bb) -- (gate);
\draw[->] (desk) -- (gate);
\draw[->,dashed] (gw.south) -- ++(0,-0.7) -| node[pos=0.25,below]{exécutions, annulations} ([xshift=8mm]gate.south);
\draw[->,dashed] (venue) -- node[right]{drop copy} (mon);
\draw[->,omThm] ([yshift=-2.5mm]gate.west) -- node[below,text=omThm]{refusé} ([yshift=-2.5mm]eng.east);
\end{tikzpicture}

\omcaption{Emplacement de la barrière~: sur le chemin de chaque ordre, avant la passerelle, alimentée par des limites assorties de \omterm{def:ll:protocols-i-fix:session}{heartbeats}, des prix de référence
et par la propre vue de la passerelle sur les exécutions et les \omterm{def:ll:order-gateway-and-order-management:inflight}{ordres en vol}. Un moniteur post-négociation branché sur le \omterm{def:ll:order-gateway-and-order-management:dropcopy}{drop copy} (chapitre~21) est une
seconde ligne, et non un substitut~: il voit ce qui s'est déjà produit.}\label{fig:ll:pre-trade-risk-and-kill-switches:where}
\end{omfigure}

\section{Contrôles en nanosecondes}

\begin{definition}[Exposition dans le pire cas]\label{def:ll:pre-trade-risk-and-kill-switches:worst}
L'\emph{exposition dans le pire cas}\index{exposition dans le pire cas} (worst-case exposure) d'une société sur un instrument est sa position, plus tout ce que ses
ordres ouverts pourraient encore ajouter s'ils étaient tous exécutés~: du côté acheteur, la position plus la quantité restante de chaque ordre d'achat ouvert,
y compris ceux dont l'accusé de réception ou l'annulation n'est pas parvenu~; du côté vendeur, la même chose avec les ordres de vente. Une limite est
respectée dans le pire cas si un nouvel ordre s'inscrit sous elle, par-dessus cette exposition.
\end{definition}

Une \omterm{def:ll:pre-trade-risk-and-kill-switches:gate}{barrière de risque} est évaluée sur chaque ordre~: elle doit donc coûter peu, et coûter autant qu'elle accepte ou qu'elle refuse~; une barrière qui est
rapide sur le chemin courant et lente sur le chemin rare est la plus lente précisément lorsqu'un emballement envoie des ordres. La barrière de la construction est
organisée pour cela. Les limites sont tenues dans des tables denses indexées par instrument (le code locate de la plateforme) et par stratégie, et non dans des
tables associatives indexées par nom. Tout ce qui peut être calculé avant l'arrivée de l'ordre l'est~: le tunnel de prix est conservé sous forme de deux prix, recalculé lorsque le
prix de référence ou les limites changent, de sorte que le contrôle du tunnel se réduit à une seule comparaison~; les expositions sont des sommes cumulées, mises à jour par
les comptes rendus comme au chapitre~21, de sorte que le contrôle de position se réduit à une addition et une comparaison. Et chaque contrôle est calculé, sans
branchement, dans un bit d'un masque (le \cref{lst:ll:pre-trade-risk-and-kill-switches:check})~; l'ordre passe si le masque vaut zéro,
et un refus signale le bit à 1 le plus bas, tandis que l'enregistrement d'audit conserve le masque entier, car la seconde raison compte lorsque
la première est corrigée.

\omcode{code/firm/riskgate/cpp/firm_riskgate.hpp}{124}{140}{Douze contrôles, un masque~: chaque condition positionne un bit sans
branchement, et un refus signale le premier.}\label{lst:ll:pre-trade-risk-and-kill-switches:check}

La \Cref{fig:ll:pre-trade-risk-and-kill-switches:cost} en donne la mesure. Avec tous les ordres acceptés et le contrôle des doublons désactivé, un contrôle
coûte environ \qty{11}{\nano\second} en médiane~; le contrôle des doublons, qui compare l'ordre aux seize plus récents de la stratégie,
ajoute environ \qty{10}{\nano\second}. La première version de ce contrôle comparait cinq champs par ordre mémorisé
et coûtait plus que tous les autres contrôles réunis~; l'empaquetage de la quantité et du prix dans un mot, et de l'instrument et du sens dans un autre
(les quantités et les prix du fil sont des champs de 32 bits, donc l'empaquetage est exact) l'a ramené à ce niveau. Sur le flux du jeu de test, où toutes les sortes de
refus se produisent et où chaque contrôle suit l'analyse de sa ligne d'entrée, il coûte environ \qty{26}{\nano\second} à la médiane,
soit le même pour les ordres acceptés et pour les ordres refusés~: le masque ne rend les refus ni moins chers ni plus coûteux.

\begin{omfigure}
\begin{tikzpicture}
\begin{axis}[width=11cm,height=5.4cm,ybar,bar width=14pt,ymin=0,ymax=45,xmin=-0.6,xmax=4.6,
  xtick={0,1,2,3,4},xticklabels={{jeu de test~:\\ tous},{jeu de test~:\\ acceptés},{jeu de test~:\\ refusés},{tous\\ acceptés},{tous acceptés~:\\ doublons\\ activés}},
  xticklabel style={align=center,font=\footnotesize},ylabel={ns par contrôle (médiane)},
  tick label style={font=\small},label style={font=\small},grid=major,grid style={gray!25},
  nodes near coords,every node near coord/.append style={font=\scriptsize},table/col sep=comma]
\addplot[fill=omThm!45,draw=omThm] table[x expr=\coordindex/2,y=ns,
  restrict expr to domain={mod(\coordindex,2)}{0:0}]{figdata/low-latency/22-pre-trade-risk-and-kill-switches/measured_check.csv};
\end{axis}
\end{tikzpicture}

\omcaption{Le coût d'un passage par les douze contrôles de la barrière, médiane des appels mesurés au TSC~: sur les 4\,874 ordres du jeu de test
(toutes les sortes de refus, tables froides), réparti selon le résultat~; et sur un flux où tous les ordres passent, avec le contrôle des doublons désactivé
puis activé. Mesures faites sur un ordinateur portable (Intel Core Ultra 7 155H) sous WSL2. Données~:
\texttt{bench\_risk.py}.}\label{fig:ll:pre-trade-risk-and-kill-switches:cost}
\end{omfigure}

Des dizaines de nanosecondes sont peu de chose à côté du trajet d'un ordre à travers le noyau et le réseau (chapitre~1), et rien ne justifie
de les réduire en contrôlant moins. La tentation va dans l'autre sens~: sortir les contrôles du processus de trading, vers un processus de risque
distinct qui voit les ordres sur un anneau et répond. Cela coûte un aller-retour entre cœurs, environ \qty{200}{\nano\second} par les
anneaux du chapitre~12, soit huit fois le contrôle lui-même, et, plus grave, cela crée une seconde défaillance à gérer~: si le processus de risque cesse de répondre, le processus de trading doit
soit attendre, soit envoyer sans contrôle. Une barrière dans le processus de trading, avec des limites qu'elle peut vérifier seule, ne tombe en panne que d'une seule façon.

\section{Limites en hiérarchie, mises à jour en direct}

Les limites se fixent à plusieurs niveaux~: la firme dispose d'un seuil de capital, chaque desk d'une part de ce seuil, chaque stratégie d'un débit, d'un nombre
d'ordres ouverts et de sa propre part, chaque instrument d'un collar, d'une taille maximale et d'une limite de position (la hiérarchie de risque du livre~6,
chapitre~29, et la politique du livre~11, chapitre~27). Un ordre est contrôlé à chaque niveau auquel il appartient, et la barrière tient à jour une
exposition courante par niveau. Les limites changent au cours de la journée, lorsqu'un trader en demande davantage ou que le desk de risque en retire,
et la barrière applique les changements de la même façon que le moteur de stratégie applique les paramètres (chapitre~20)~: par instantanés complets et versionnés,
entre deux ordres, jamais au milieu d'un contrôle, afin qu'aucun ordre ne soit contrôlé en partie selon un ancien instantané et en partie
selon un nouveau.

La barrière se ferme en cas de défaillance. Un instantané porte un âge maximal~: le service de risque envoie des \omterm{def:ll:protocols-i-fix:session}{heartbeats}, et si le dernier
instantané ou \omterm{def:ll:protocols-i-fix:session}{heartbeat} est plus vieux que cet âge, chaque ordre est refusé comme périmé. Le jeu de test de la construction comprend un service de risque
qui se tait pendant trois secondes avec un âge maximal de deux~: dès que son dernier \omterm{def:ll:protocols-i-fix:session}{heartbeat} a deux secondes, et jusqu'à l'arrivée du suivant,
la barrière refuse chaque ordre, soit 50 ordres. Il en va de même pour les prix de référence qui servent au collar~: un instrument dont la référence est
plus vieille que son âge maximal ne peut pas être négocié, car un collar autour d'un prix périmé n'est pas un collar. Se fermer en cas de défaillance a un coût, puisqu'une stratégie
s'arrête lorsqu'un service de surveillance tombe en panne, et la firme doit décider de quels services dépend la négociation. La règle de la construction est que
tout ce dont un contrôle a besoin est un tel service.

\section{Kill switches~: bloquer, annuler, aplatir}

\begin{definition}[Annulation de masse]\label{def:ll:pre-trade-risk-and-kill-switches:mass}
Une \emph{annulation de masse}\index{annulation de masse} (mass cancel) est une requête qui annule beaucoup d'ordres à la fois, tous les ordres ouverts d'une session,
d'un instrument ou d'un sens, soit comme un seul message que la plateforme offre (le message \texttt{M} du simulateur), soit comme une annulation pour chaque
ordre ouvert, envoyées ensemble~; les ordres restent en vol, et peuvent encore être exécutés, jusqu'à ce que chaque annulation soit confirmée.
\end{definition}

Un kill switch arrête un nœud de la hiérarchie (une stratégie, un desk, la firme) et peut faire trois choses, par ordre croissant de gravité.
\emph{Bloquer} refuse chaque nouvel ordre sous le nœud~: rien de plus ne sort, et rien de ce qui est déjà sorti n'est touché.
\emph{Annuler} envoie en plus une \omterm{def:ll:pre-trade-risk-and-kill-switches:mass}{annulation de masse} pour les ordres ouverts du nœud~; tant que chaque annulation n'est pas confirmée, les ordres comptent dans
l'exposition, et ceux qui sont exécutés dans l'intervalle sont des situations de course du type de celles du chapitre~21. \emph{Aplatir} envoie en plus des ordres qui
clôturent les positions, au prix de référence, que le collar admet~; c'est le plus dangereux des trois, car il négocie
dans un marché que l'emballement a pu déplacer, et la construction renvoie ces ordres à une personne pour confirmation au lieu de les envoyer.
Qui actionne le kill switch relève de la politique (livre~11)~: la barrière le déclenche lors d'une violation d'une limite stricte, le desk de risque peut
le déclencher à la main, et beaucoup de plateformes en proposent un propre. Le réarmement relève toujours d'une décision humaine.

\omcode{code/firm/riskgate/firm_riskgate.py}{198}{210}{Le kill switch de la référence Python~: bloquer, annuler (une annulation de masse
des ordres ouverts du nœud), aplatir (les ordres de clôture, renvoyés pour confirmation).}\label{lst:ll:pre-trade-risk-and-kill-switches:kill}

\section{Où se placent les contrôles}

Les contrôles peuvent se placer à quatre endroits, et une firme en utilise plusieurs. Dans le processus de trading, comme ici~: l'endroit le moins coûteux et le seul
qui voie chaque ordre avant qu'il ne parte, avec la vue propre de la firme sur les \omterm{def:ll:order-gateway-and-order-management:inflight}{ordres en vol}. Dans un processus de risque distinct ou un appareil
placé sur le trajet~: indépendant du code de trading, au prix d'un saut supplémentaire, et la seule protection contre un bogue dans
le processus de trading lui-même. Chez le courtier, lorsque la firme atteint la plateforme par l'accès sponsorisé d'un courtier
(livre~1, chapitre~4)~: le courtier répond, au titre de la Market Access Rule, des ordres qu'il laisse passer, quels que soient les
contrôles de la firme. Et à la plateforme, dont les collars, les limites et les kill switches protègent le marché plutôt que la firme. Le routeur de 2012
avait des contrôles en amont et une surveillance en aval, et rien sur le trajet de sa propre sortie. Un test utile de toute architecture consiste à
demander, pour chaque composant qui génère des ordres, quel contrôle voit sa sortie avant que la plateforme ne la voie.

\section{Tutoriel~: une course, un contrôle à la fois}\label{tut:ll:pre-trade-risk-and-kill-switches:tut}

\texttt{firm\_riskgate\_runaway.py} reconstruit la forme de 2012 face au moteur d'appariement du livre~10. La course envoie des ordres enfants d'achat
de 100 actions par rafales de quinze, espacés d'une microseconde, toutes les \qty{10}{\milli\second}~: 1\,500 par seconde en moyenne, soit le rythme de
4 millions d'exécutions en 45 minutes, d'environ 99 actions chacune. Elle les cote cinq ticks au-dessus de la dernière transaction et ne regarde jamais
ses exécutions. Un fournisseur de liquidité tient une échelle d'ordres de vente et la remonte d'un tick toutes les dix exécutions, de sorte que la course fait monter
le prix à mesure qu'elle achète. Le \Cref{tab:ll:pre-trade-risk-and-kill-switches:runaway} et
la \cref{fig:ll:pre-trade-risk-and-kill-switches:runaway} en donnent cinq secondes, chaque contrôle étant activé à tour de rôle.

\begin{table}[ht]
\centering\small
\begin{tabular}{@{}>{\raggedright\arraybackslash}p{5.6cm}>{\centering\arraybackslash}p{2.3cm}>{\centering\arraybackslash}p{2.3cm}>{\centering\arraybackslash}p{2.3cm}>{\centering\arraybackslash}p{2.3cm}@{}}
\toprule
contrôle activé & ordres envoyés & actions achetées & premier refus & reporté sur 45 min \\
\midrule
aucun & 7\,500 & 750\,000 & --- & 405 millions \\
collar, 5~\% autour de la dernière transaction & 7\,500 & 750\,000 & --- & 405 millions \\
taille et notionnel par ordre & 7\,500 & 750\,000 & --- & 405 millions \\
limitation, 500 ordres par seconde & 2\,545 & 254\,500 & \qty{40}{\milli\second} & 135 millions \\
doublons dans \qty{100}{\milli\second} & 55 & 5\,500 & première répétition & 3.0~millions \\
position 20\,000, exécutions seulement & 210 & 21\,000 & \qty{140}{\milli\second} & 21\,000 \\
position 20\,000, en vol & 200 & 20\,000 & \qty{130}{\milli\second} & 20\,000 \\
seuil de capital, 2.5~millions de dollars & 249 & 24\,900 & \qty{160}{\milli\second} & 24\,900 \\
\bottomrule
\end{tabular}
\caption{Une course de la forme de 2012, cinq secondes face au moteur d'appariement, avec un contrôle à la fois. «~Reporté sur
45 minutes~» prolonge le rythme de la dernière seconde lorsque la course envoie encore, et conserve la position lorsqu'elle a été
arrêtée. Données~: \texttt{fig\_runaway.py} (déterministe).}\label{tab:ll:pre-trade-risk-and-kill-switches:runaway}
\end{table}

\begin{omfigure}
\begin{tikzpicture}
\begin{axis}[width=11cm,height=6cm,ymode=log,xmin=0,xmax=5000,ymin=50,ymax=2e6,
  xlabel={temps depuis le début de la course (ms)},ylabel={position (actions)},
  tick label style={font=\small},label style={font=\small},grid=major,grid style={gray!25},
  legend pos=outer north east,legend style={font=\scriptsize,cells={anchor=west}},table/col sep=comma]
\addplot[omThm,thick] table[x=t_ms,y=none]{figdata/low-latency/22-pre-trade-risk-and-kill-switches/runaway_trace.csv};
\addlegendentry{aucun (ou collar, taille)}
\addplot[omMeth,thick] table[x=t_ms,y=throttle]{figdata/low-latency/22-pre-trade-risk-and-kill-switches/runaway_trace.csv};
\addlegendentry{limitation}
\addplot[gray,thick,dashed] table[x=t_ms,y=duplicates]{figdata/low-latency/22-pre-trade-risk-and-kill-switches/runaway_trace.csv};
\addlegendentry{doublons}
\addplot[omDef,thick] table[x=t_ms,y=position_in_flight]{figdata/low-latency/22-pre-trade-risk-and-kill-switches/runaway_trace.csv};
\addlegendentry{position, en vol}
\addplot[black,thick,dotted] table[x=t_ms,y=capital_threshold]{figdata/low-latency/22-pre-trade-risk-and-kill-switches/runaway_trace.csv};
\addlegendentry{seuil de capital}
\end{axis}
\end{tikzpicture}

\omcaption{La position du fuyard, toutes les \qty{50}{\milli\second}, avec un seul contrôle à la fois. Les contrôles portant sur chaque ordre isolément
ne se déclenchent jamais~; un débit limité ralentit le fuyard et le contrôle des doublons le freine presque entièrement, tandis qu'une limite de position qui compte les \omterm{def:ll:order-gateway-and-order-management:inflight}{ordres en vol},
ou un seuil de capital, l'arrête en moins de 250 ordres. Données~: \texttt{fig\_runaway.py}.}\label{fig:ll:pre-trade-risk-and-kill-switches:runaway}
\end{omfigure}

Le tableau constitue l'argument du chapitre. Les contrôles qui regardent chaque ordre seul (collar, taille, notionnel) ne se déclenchent jamais~: chaque ordre enfant
est petit et fixé près d'un marché que le fuyard fait lui-même évoluer, si bien qu'un collar autour de la dernière transaction se déplace avec lui
(le cours a monté de 7.5~\% en cinq secondes, davantage que le collar, et le collar ne l'a jamais remarqué). Un débit limité ralentit le fuyard jusqu'à
son rythme et ne l'arrête pas~: 135 millions d'actions en 45 minutes restent une catastrophe. Le contrôle des doublons freine presque ce
fuyard, parce que ses ordres enfants se répètent, mais raterait celui qui varierait sa quantité. Ce qui l'arrête est une limite sur ce que la
firme détient, comptée dans le pire cas~: la limite de position à 20\,000 actions après 200 ordres, le seuil de capital à
24\,900. Si l'on ne compte que les exécutions, le dépassement porte sur les \omterm{def:ll:order-gateway-and-order-management:inflight}{ordres en vol}~: ici une rafale de dix ordres, soit 1\,000 actions~; avec davantage d'\omterm{def:ll:order-gateway-and-order-management:inflight}{ordres en vol},
ou un chemin de rapport plus lent, le dépassement est plus grand. Le kill switch change ce qui se passe ensuite~: sans lui, le fuyard
continue d'envoyer et la barrière continue de refuser, soit 7\,300 refus en cinq secondes~; avec lui, la stratégie est bloquée à la première
violation et ses ordres ouverts sont annulés.

\section{Construction~: la barrière de risque}\label{bld:ll:pre-trade-risk-and-kill-switches:riskgate}

\begin{build}
\bfield{Rôle} Le contrôle sur le chemin de chaque ordre, entre le moteur de stratégie (chapitre~20) et la \omterm{def:ll:order-gateway-and-order-management:gateway}{passerelle d'ordres}
(chapitre~21)~; la politique qu'il applique est celle du livre~11 (\texttt{firm.riskctl}, chapitre~27), et ses enregistrements d'audit vont au
journal binaire (chapitre~23).
\bfield{Interface} Référence Python \texttt{firm\_riskgate}~: \texttt{Limits.from\_dict}, \texttt{from\_riskctl},
\texttt{Gate(limits, t)} avec \texttt{check}, \texttt{set\_limits}, \texttt{heartbeat}, \texttt{set\_reference},
\texttt{on\_fill}, \texttt{on\_done}, \texttt{kill}, \texttt{unkill}, \texttt{audit}~; \texttt{replay}, \texttt{decisions}~;
\texttt{firm\_riskgate\_runaway.run} contre le moteur du livre~10. C++20 \texttt{firm::risk}~: \texttt{parse\_limits},
\texttt{Gate} avec les mêmes opérations et \texttt{Decision\{code, mask\}}. Rust \texttt{firm\_riskgate}~: \texttt{Gate},
\texttt{parse\_limits}, \texttt{replay}.
\bfield{Règles} Chaque contrôle est calculé dans un masque, sans branchement sur le résultat~; les limites sont des instantanés complets entre deux ordres~; les limites
ou les prix de référence périmés entraînent un refus (à défaillance fermée)~; l'exposition compte chaque \omterm{def:ll:order-gateway-and-order-management:inflight}{ordre en vol}~; aucune allocation par contrôle~; un kill ne se lève
que par une personne.
\bfield{Tests d'acceptation} \texttt{code/firm/riskgate/}~: les 4\,874 ordres du jeu de test rejoués aux mêmes décisions en Python, C++
et Rust, chaque contrôle refusant au moins une fois~; chaque contrôle à sa frontière~; défaillance fermée sur les limites et les références périmées~; les trois
actions du kill switch au niveau de la firme, du desk et de la stratégie~; les sommes courantes égales à un recomptage après des flux aléatoires~; zéro
allocation~; le fuyard arrêté par la limite de position et le seuil de capital, et non par le collar~; les limites \texttt{firm.riskctl} du livre~11,
exportées pour la barrière, rejouées en Python et C++.
\bfield{Pour aller plus loin} Une limite de perte sur le profit réalisé et non réalisé (exercice~7)~; des limites de messages par plateforme~; des limites partagées
par plusieurs passerelles via une mémoire partagée, avec un seul écrivain~; les moniteurs du livre~11 (perte, débit de messages) alimentant le kill switch.
\end{build}

\begin{omsources}
\item U.S. Securities and Exchange Commission, Release No. 34-70694, \emph{In the Matter of Knight Capital Americas LLC}, 2013.
\item 17 CFR 240.15c3-5, \emph{Risk management controls for brokers or dealers with market access}.
\item Commission Delegated Regulation (EU) 2017/589 of 19 July 2016 (RTS 6).
\end{omsources}

\section{Exercices}

\begin{exercise}[$\star$]\label{exo:ll:pre-trade-risk-and-kill-switches:1}
Le prix de référence est 25.0000 (250\,000 en unités de \num{e-4}) et le collar est de 500 points de base. Lesquels des ordres d'achat à
26.2500 et 26.2600, et des ordres de vente à 23.7500 et 23.7400, passent~?
\end{exercise}

\begin{exercise}[$\star$]\label{exo:ll:pre-trade-risk-and-kill-switches:2}
La limite longue est de 5\,000 actions. La position est de 2\,000 longues, et des ordres d'achat de 1\,500 et 1\,000 sont ouverts, dont l'un n'est pas
encore acquitté. Un achat de 600 est-il accepté~? Quel est le plus grand achat qui l'est~?
\end{exercise}

\begin{exercise}[$\star$]\label{exo:ll:pre-trade-risk-and-kill-switches:3}
À partir des totaux du régulateur, calculer le nombre moyen d'exécutions par seconde et d'actions par exécution de l'incident.
\end{exercise}

\begin{exercise}[$\star\star$]\label{exo:ll:pre-trade-risk-and-kill-switches:4}
Une rafale de quinze ordres enfants identiques part à une microseconde d'intervalle. La fenêtre de doublons est de \qty{100}{\milli\second}. Combien
de ces quinze ordres passent, et comment le fuyard devrait-il changer pour tous les laisser passer~?
\end{exercise}

\begin{exercise}[$\star\star$]\label{exo:ll:pre-trade-risk-and-kill-switches:5}
L'âge maximal des limites est de deux secondes et les \omterm{def:ll:protocols-i-fix:session}{heartbeats} arrivent toutes les \qty{400}{\milli\second}. Le service de risque s'arrête à l'instant
$t$. Jusqu'à quand les ordres sont-ils acceptés, et pourquoi «~jusqu'au retour du service~» est-ce la mauvaise réponse~?
\end{exercise}

\begin{exercise}[$\star\star$]\label{exo:ll:pre-trade-risk-and-kill-switches:6}
Un contrôle coûte environ \qty{26}{\nano\second} sur le flux du jeu de test. Quelle fraction d'un budget tick-to-trade de
\qty{1}{\micro\second} cela représente-t-il, et que coûterait en plus un contrôle dans un processus séparé sur un autre cœur~?
\end{exercise}

\begin{exercise}[$\star\star\star$]\label{exo:ll:pre-trade-risk-and-kill-switches:7}
\emph{Programmation.} Ajouter une limite de perte par stratégie~: le profit réalisé issu des exécutions plus le profit non réalisé au prix de référence, avec
un kill (blocage et annulation) lorsqu'il descend sous la limite. Où la perte est-elle mise à jour, et pourquoi n'est-elle pas un contrôle pré-négociation comme les
autres~?
\end{exercise}

\begin{exercise}[$\star\star\star$]\label{exo:ll:pre-trade-risk-and-kill-switches:8}
\emph{Trouver la faille.} «~Nos contrôles de risque tournent dans leur propre processus, qui lit une copie de chaque ordre envoyé par la passerelle et envoie
un kill dès qu'il voit une violation. Ils restent ainsi hors du chemin de latence.~»
\end{exercise}

\section{Problème~: quarante-cinq minutes}

\begin{problem}[{Problème du week-end --- quels contrôles auraient arrêté le fuyard de 2012}]\label{pb:ll:pre-trade-risk-and-kill-switches:1}
Utiliser le \cref{tab:ll:pre-trade-risk-and-kill-switches:runaway} (\texttt{runaway.csv}) et les totaux du régulateur~: 4 millions
d'exécutions, plus de 397 millions d'actions, environ 45 minutes.

\medskip\noindent\textbf{Partie I --- Les débits.}
\begin{enumerate}
\item Calculer le débit d'exécution de l'incident et vérifier que les 1\,500 ordres par seconde du fuyard lui correspondent.
\item Combien d'actions le fuyard non contrôlé achète-t-il en cinq secondes, et en 45 minutes au même débit~? Comparer avec
l'incident.
\item Dans la simulation, le prix a monté d'un tick toutes les dix exécutions. De combien a-t-il monté en cinq secondes~?
\item Pourquoi le collar ne s'est-il jamais déclenché, alors que le prix a bougé de plus que le collar~?
\end{enumerate}

\medskip\noindent\textbf{Partie II --- Les contrôles qui ne l'arrêtent pas.}
\begin{enumerate}[resume]
\item Pourquoi les limites de taille et de notionnel par ordre ne se déclenchent-elles jamais~?
\item Le débit limité se déclenche après \qty{40}{\milli\second}. Pourquoi si tard, et quel est ensuite le débit du fuyard~?
\item Reporter le débit limité sur 45 minutes.
\item Le contrôle des doublons laisse passer 55 ordres en cinq secondes. Lesquels, et quel petit changement du fuyard le mettrait en
échec~?
\end{enumerate}

\medskip\noindent\textbf{Partie III --- Les contrôles qui l'arrêtent.}
\begin{enumerate}[resume]
\item Combien d'ordres la limite de position de 20\,000 actions, comptée en vol, laisse-t-elle passer, et à partir de quel moment se déclenche-t-elle~?
\item La limite qui ne compte que les exécutions s'arrête à 21\,000. Expliquer la différence.
\item Dans quelles circonstances cette différence serait-elle importante~?
\item Le seuil de capital est de 2.5~millions de dollars. À environ 100~\$ par action, pourquoi arrête-t-il le fuyard vers 24\,900 actions~?
\end{enumerate}

\medskip\noindent\textbf{Partie IV --- Le verdict.}
\begin{enumerate}[resume]
\item Énoncer le \emph{résultat nommé}~: ordres et actions avant le déclenchement de chaque contrôle, reportés sur 45 minutes, face aux
totaux de l'incident.
\item Que le kill switch ajoute-t-il à la limite de position dans cette simulation~?
\item Le fuyard a négocié 154 titres. Qu'est-ce que cela change pour une limite de position par instrument, et quel contrôle y est insensible~?
\item Lequel des contrôles manquants listés dans l'ordonnance du régulateur correspond à chaque ligne du tableau~?
\item Pourquoi la barrière doit-elle être à défaillance fermée~?
\item Pourquoi un moniteur post-négociation, aussi rapide soit-il, ne remplace-t-il pas la barrière~?
\item Quels faux positifs le contrôle des doublons déclencherait-il un jour ordinaire~?
\item En une phrase~: quel type de limite arrête un fuyard~?
\end{enumerate}
\end{problem}

\section{Questions d'entretien}

\begin{interviewq}[$\star$ \iqroles{développeur}]\label{iq:ll:pre-trade-risk-and-kill-switches:1}
Listez les contrôles pré-négociation que vous placeriez sur chaque ordre, et dites lesquels vous attendriez pour repérer un algorithme fuyard.
\end{interviewq}

\begin{interviewq}[$\star\star$ \iqroles{développeur}]\label{iq:ll:pre-trade-risk-and-kill-switches:2}
Comment faites-vous pour qu'un contrôle de risque coûte quelques nanosecondes, et qu'il coûte autant lorsqu'il accepte que lorsqu'il refuse~?
\end{interviewq}

\begin{interviewq}[$\star\star$ \iqroles{développeur, risque}]\label{iq:ll:pre-trade-risk-and-kill-switches:3}
Votre service de risque cesse d'envoyer ses limites. Que fait le système de négociation, et pourquoi~?
\end{interviewq}

\begin{interviewq}[$\star\star$ \iqroles{développeur, trader}]\label{iq:ll:pre-trade-risk-and-kill-switches:4}
Que doit faire un kill switch, et qui doit pouvoir le déclencher et l'annuler~?
\end{interviewq}

\begin{interviewq}[$\star\star$ \iqroles{développeur}]\label{iq:ll:pre-trade-risk-and-kill-switches:5}
Comment modifiez-vous les limites d'une stratégie en cours de journée sans l'arrêter~?
\end{interviewq}

\begin{interviewq}[$\star\star\star$ \iqroles{développeur, risque}]\label{iq:ll:pre-trade-risk-and-kill-switches:6}
Lire l'incident de 2012 comme un document d'exigences. Concevez les contrôles, leur emplacement, et comment vous vérifieriez qu'ils arrêtent un
fuyard.
\end{interviewq}
